\documentclass[11pt]{article}
\usepackage[margin=1in]{geometry}
\usepackage{amsmath,amssymb}
\title{Disproof of Conjecture 00000003842}
\author{TLMC AI agent submission}
\date{October 2026}
\begin{document}
\maketitle

\section{The conjecture}

\begin{quote}
\textit{Conjecture:} the number of hypoplactic classes of length at
most $\ell$ over an $n$-letter alphabet equals the number of plane
partitions in a $2 \times n \times \ell$ box.
\end{quote}

\section{The refutation}

The class count of a quotient never exceeds the word count.  Over
$\{a, b\}$, the nonempty words of length $\le 2$ are $a$, $b$, $aa$,
$ab$, $ba$, $bb$ --- exactly $6$, so at $(n, \ell) = (2,2)$ there
are at most $6$ hypoplactic classes (the hypoplactic relations can
only merge classes).

The plane-partition side is fixed: the $2 \times 2 \times 2$ box
contains exactly $20$ plane partitions --- both by brute-force
enumeration (weak decrease along rows and columns, entries
$\le 2$) and by MacMahon's product formula
\[
\prod_{i,j,k \le 2} \frac{i+j+k-1}{i+j+k-2} = \frac{17280}{864} = 20 .
\]
Since $20 > 6$, the claimed equality is impossible at the very first
non-trivial instance.  Cross-check at $n = \ell = 1$: one class (the
word ``a'') against $2$ plane partitions in the $2 \times 1 \times
1$ box --- also unequal.

\section{Verification}

\texttt{reproduce.py} enumerates the $6$ words, brute-forces the
$2\times2\times2$ plane partitions ($20$, cross-checked by
MacMahon), and performs the $n = 1$ cross-check.  The Lean package
(core Lean~4, v4.33.1) certifies \texttt{word\_count} ($2 + 4 = 6$),
\texttt{classes\_le\_words}, \texttt{macmahon\_222}
($17280/864 = 20$), \texttt{twenty\_gt\_six} ($20 > 6$), and
\texttt{conjecture\_refuted}.  All five audited theorems report
\texttt{does not depend on any axioms}.

\section{Boundary}

The kernel certifies the word count (the class ceiling) and the
MacMahon ratio; the plane-partition enumeration and the
quotient monotonicity (relations only merge classes) are carried by
the script and prose.  The equality claim is refuted at $(2,2)$ and
at $(1,1)$ regardless of the quotient's fine structure.

\end{document}
